\chapter{Proposed Solution}
\label{ch:proposed-solution}

\section{Proposed Solution}
\label{sec:proposed-solution}

The proposed solution is a mobile application that supports children with speech impairments through guided and interactive speech practice. The application is designed as a supplementary tool that can help structure practice outside traditional therapy sessions.

Describe the main components of the solution and explain why each part addresses the problem introduced in \Cref{ch:problem-statement}.

Possible components:
\begin{itemize}[noitemsep]
  \item Guided speech exercises.
  \item Child-friendly interface and interaction design.
  \item Audio or speech input.
  \item Feedback after exercises.
  \item Progress tracking.
  \item Caregiver or therapist support features.
  \item Backend services and data storage, if applicable.
\end{itemize}

\section{Engineering Rationale}
\label{sec:engineering-rationale}

Explain why the chosen solution is a good fit from a software engineering perspective. Discuss criteria such as accessibility, usability, maintainability, scalability, privacy, and technical feasibility.

\begin{table}[H]
  \centering
  \caption{Example table for comparing solution design considerations.}
  \label{tab:solution-considerations}
  \begin{tabular}{p{0.28\textwidth} p{0.58\textwidth}}
    \toprule
    \textbf{Consideration} & \textbf{Relevance to the project} \\
    \midrule
    Accessibility & Mobile devices can make speech practice available outside scheduled therapy sessions. \\
    Usability & A child-friendly interface can reduce friction and support repeated use. \\
    Maintainability & A modular implementation makes the system easier to extend and test. \\
    Privacy & Speech-related data and child user data require careful handling. \\
    \bottomrule
  \end{tabular}
\end{table}
